\chapter{Ewaluacja}\label{ch:ewaluacja}

% TODO: wyniki ewaluacji — ok. 8--15 stron.

\section{Metodyka ewaluacji}\label{sec:metodyka-ewaluacji}

% TODO: opis zbioru testowego, metryki, narzędzia pomiarowe.

\section{Zbiór testowy (ground truth)}\label{sec:ground-truth}

% TODO: opis scenariuszy testowych, sposób ich anotacji, liczba próbek
% per specjalizacja.

\section{Jakość transkrypcji (STT)}\label{sec:ewaluacja-stt}

% TODO: metryki WER i CER, porównanie konfiguracji faster-whisper,
% wpływ diaryzacji na jakość.

\begin{table}[htbp]
  \centering
  \caption{Przykładowe wyniki WER dla różnych konfiguracji STT
           (dane zastępcze --- do uzupełnienia).}
  \label{tab:wer}
  \begin{tabular}{lrrr}
    \toprule
    Konfiguracja & WER [\%] & CER [\%] & RTF \\
    \midrule
    whisper-small       & --- & --- & --- \\
    whisper-medium      & --- & --- & --- \\
    whisper-large-v3    & --- & --- & --- \\
    \bottomrule
  \end{tabular}
\end{table}

\section{Jakość strukturyzacji SOAP}\label{sec:ewaluacja-slm}

% TODO: ewaluacja eksperta klinicznego, metryki spójności, kompletności,
% poprawności medycznej.

\section{Skuteczność anonimizacji}\label{sec:ewaluacja-anonimizacja}

% TODO: precyzja/czułość rozpoznawania PII, współczynnik wycieku danych,
% analiza błędów.

\section{Trafność sugestii RAG i NFZ}\label{sec:ewaluacja-rag}

% TODO: metryki Recall@k, Precision@k, NDCG; akceptowalność sugestii
% przez lekarzy w~scenariuszach testowych.

\section{Wydajność systemu end-to-end}\label{sec:ewaluacja-wydajnosc}

% TODO: czas odpowiedzi per etap, wykorzystanie zasobów GPU/CPU,
% skalowanie przy większej liczbie równoległych wizyt.

\section{Dyskusja wyników}\label{sec:dyskusja}

% TODO: interpretacja wyników, porównanie z~hipotezami, identyfikacja
% ograniczeń.
